\documentclass{ltjsbook}
\usepackage{docmute} 
\usepackage{../../myStyle} 

\begin{document}
\chapter{行列による関係の表現}
\label{x07_matrixModel}
前回から線形代数の話をしています。線形代数は数字をセットで扱うための表現方法，計算方法ですから，多変量データを分析しようという時には必須の技術になります。
前回は線形代数の導入ですから，計算方法を解説してきましたが，今回はこの計算方法を使って具体的にデータをどのように表現し，どのように計算するのかを見ていくことになります。

\section{データの行列表現}

ここまで行列の形ばかり見て来ましたが，狙いはあくまでも調査研究など，多変量データを扱う場面での利用です。なぜ多変量データ分析をする際にこのような知識が必要なのか，思うかもしれません。ですが，得られるデータは行列として扱うと表現が大変便利なのです。たとえば質問項目が$m$個あって，調査対象者$n$人から回答を得たとすると，データは次のように表現できます。

\[
	\bm{X}=
	\begin{pmatrix}
		x_{11} & x_{12} & \cdots & x_{1m} \\
		x_{21} & x_{22} & \cdots & x_{2m} \\
		\vdots & \vdots & \ddots & \vdots \\
		x_{n1} & x_{n2} & \cdots & x_{nm} \\
	\end{pmatrix}
\]

データ全体をこうして，ひとつの記号で表現できたら便利ですよね。これらを使ったデータの表記に慣れておきましょう。

各反応の平均点は以下のように表現されます。まず，要素がすべて$1$からなるベクトルを次のように表します。
\[ \bm{1}=\begin{pmatrix} 1 \\ 1 \\ \vdots \\1 \end{pmatrix} \]
わかりにくいかもしれませんが，この$\bm{1}$は太字でベクトルを表しており，スカラーの$1$とは違うことに注意してください。

さて，各項目の和はベクトルの掛け算の定義によって次のように表現できます。
\[ \bm{X}' \bm{1} = \begin{pmatrix} \sum_{i=1}^{n} x_{i1} \\ \sum_{i=1}^{n} x_{i2} \\ \vdots \\ \sum_{i=1}^{n} x_{im} \end{pmatrix} \]


これを使って平均値(列)ベクトル$\bm{m}$を次のように表すことができます。
\[ \bm{m} = \frac{1}{n} \bm{X}' \bm{1} =
	\begin{pmatrix} 1/n\sum_{i=1}^{n} x_{i1} \\ 1/n\sum_{i=1}^{n} x_{i2} \\ \vdots \\ 1/n\sum_{i=1}^{n} x_{im} \end{pmatrix} =
	\begin{pmatrix}\bar{x_.1} \\ \bar{x_.2}\\ \vdots \\\bar{x_.m} \end{pmatrix} \]

ここで$\bar{x_.1}$は1つめの添字$i$を足し合わせて割ることでなくしていますから，$\bar{x_1}$のように省略して書くことがあります。このときの$1$は「第一番目の変数」という意味であり，個人の情報がなくなっている変数を意味していることに注意してください。

さて，平均からの偏差を要素に持つ行列$\bm{V}$を考えたとします。
\[
	\begin{array}{ccl}
		\bm{V} & = & \bm{X} - \bm{1} \bm{m}'                                                                                                                  \\
		       &   &                                                                                                                                          \\
		       & = & \bm{X} - \begin{pmatrix} 1 \\ 1 \\ \vdots \\ 1 \end{pmatrix} \begin{pmatrix} \bar{x_.1} & \bar{x_.2} & \cdots & \bar{x_.m} \end{pmatrix} \\
		       &   &                                                                                                                                          \\
		       & = & \bm{X} - \begin{pmatrix} \bar{x_.1} & \bar{x_.2} & \cdots & \bar{x_.m} \\
                \bar{x_.1} & \bar{x_.2} & \cdots & \bar{x_.m} \\
                \vdots     & \vdots     & \ddots & \vdots     \\
                \bar{x_.1} & \bar{x_.2} & \cdots & \bar{x_.m}\end{pmatrix}
	\end{array}
\]
この行列$\bm{V}$のサイズは$n \times m$であることに注意してください。
これはまた，次のように表すこともできます。
\[ \bm{V} = ( \bm{I} -\frac{1}{n}\bm{1} \bm{1} ')\bm{X} \]
ここで$\bm{I}$は適切なサイズの単位行列です\footnote{適切なサイズってなんだよ，と思いますよね。これは計算に合うようなサイズ，という意味です。具体的に考えてみますと，$\bm{I}$の後ろは$\bm{1}\bm{1}'$です。$\bm{1}$は$n\times 1$の列ベクトルで，転置したものと掛け合わせますから，$\bm{1}\bm{1}'$のサイズは$n \times n$です。行列の引き算は同じサイズでないと成立しませんから，ここでの$\bm{I}$も$n \times n$でなければなりません。カッコの中身が$n \times n$で，それにサイズ$n \times m$である$\bm{X}$をかけますから，計算結果や右辺のサイズは$n \times m$になります。}。

これを使うと，たとえば分散共分散行列$\bm{S}$は次のようになります。
\[ \bm{S} = \frac{1}{n} \bm{V}' \bm{V} = \begin{pmatrix}
		s_1^2  & s_{12} & \cdots & s_{1m} \\
		s_{21} & s_2^2  & \cdots & s_{2m} \\
		\vdots & \vdots & \ddots & \vdots \\
		s_{m1} & s_{m2} & \cdots & s_m^2\end{pmatrix} \]
ここで$s_j$とあるのは第$j$変数の標準偏差を，$s_{jk}$とあるのは第$j$変数と第$k$変数の共分散です。添え字は変数番号になっています。
また，ここでもサイズに注目してください。$\bm{V}'\bm{V}$は，サイズで言うと$m \times n$と$n \times m$の積ですから，$m \times m$ になります。この行列は正方対称行列です。

また，対角項に各変数の標準偏差$s_j$が入った行列$\bm{Q}$を以下のように定めるとしましょう。次のような行列です。
\[ \bm{Q} = \begin{pmatrix} s_1    & 0      & \cdots & 0      \\
                0      & s_2    & \cdots & 0      \\
                \vdots & \vdots & \ddots & \vdots \\
                0      & 0      & \cdots & s_m\end{pmatrix} \]
そうすると，これの\keytermL{逆行列}{ぎゃくぎょうれつ}をつかって標準得点行列$\bm{Z}$を次のように表すことができます。
\[ \bm{Z} = \bm{V} \bm{Q}^{-1} \]
さらに，これを用いて相関行列$\bm{R}$を次のように表すことができます。
\[ \bm{R} = \frac{1}{n}\bm{Z}' \bm{Z} = \begin{pmatrix} 1      & r_{12} & \cdots & r_{1m} \\
                r_{21} & 1      & \cdots & r_{2m} \\
                \vdots & \vdots & \ddots & \vdots \\
                r_{m1} & r_{m2} & \cdots & 1\end{pmatrix} \]
データのサイズにかかわらず，一般的にこのように表現できるのはとてもわかりやすいですね。

\section{線形モデルの行列表現}
ベクトルや行列の記法をつかうと，\keytermL{回帰分析}{かいきぶんせき}や\keytermL{重回帰分析}{じゅうかいきぶんせき}の式がとても単純な形で表現できます。

回帰分析は，$Y=aX+b+e$という式で表現できる，ということでしたが，式中の$X$や$Y$は観測されたデータですので，$X=\{x_1,x_2,\cdots,x_n\}$，$Y=\{y_1,y_2,\cdots,y_n\}$というベクトルだと考えることができます。ですから，正確に書けば，ベクトルを使って次のように書くべきです。
\[ \bm{Y} = a\bm{X} + \bm{b} + \bm{e}\]
これは，要素を表現しながら書くと\footnote{エレメントワイズelement-wiseの表現，と言ったりします。}次のようになります。
\[
	\begin{pmatrix}
		y_1    \\
		y_2    \\
		\vdots \\
		y_n
	\end{pmatrix}
	=a
	\begin{pmatrix}
		x_1    \\
		x_2    \\
		\vdots \\
		x_n
	\end{pmatrix}
	+
	\begin{pmatrix}
		b      \\
		b      \\
		\vdots \\
		b
	\end{pmatrix}
	+
	\begin{pmatrix}
		e_1    \\
		e_2    \\
		\vdots \\
		e_n
	\end{pmatrix}
\]

列ベクトル$\bm{X}$に定数$a$をかけて得られるのは同じサイズの列ベクトル，列ベクトル同士は足しても同じサイズの列ベクトルですから，左辺と右辺はどちらも列ベクトルで，対応関係が取れていることになります。

ここで少し表現の工夫をします。説明変数$X$のベクトルの左に数字の1だけが入った列を作ります。また，係数もまとめてベクトル$\bm{\beta}$を次のように用意します。
\[
	\bm{X} =
	\begin{pmatrix}
		1      & x_1    \\
		1      & x_2    \\
		\vdots & \vdots \\
		1      & x_n
	\end{pmatrix},
	\bm{\beta}=
	\begin{pmatrix}
		b \\
		a
	\end{pmatrix}
\]
このようにすると，回帰分析の式は
\[ \bm{Y} = \bm{X\beta} + \bm{e} \]
と表すことができます。とても簡単な表現になりました(試しに各行を行列の計算式に則って計算してみてください。うまく表現できていることがわかると思います)。

さらにこの表現はありがたいことに，複数の説明変数がある重回帰分析の時でも同じ形で表すことができます。重回帰分析は，これまでの書き方ですと$Y = a_1 X_1 + a_2 X_2 + \cdots + a_n X_n + b + e$というようにしていました。ここで，係数と切片を1つの行列で表現する時わかりやすくするために，少し書き換えて$Y = \beta_0 + \beta_1 X_1 + \beta_2 X_2 +\cdots +e $としましょう。記号が変わっただけで中身は同じ，意味も同じです\footnote{厳密に言えば記述統計学として誤差を最小にするように推定した係数はアルファベット$b_0,b_1,...$で，推測統計学として母数の推定値として算出した係数はギリシア文字$\beta_0,\beta_1....$で表現する，というルールです。すでに習ったように，最小二乗法での推定値と最尤法での推定値は，誤差が正規分布する場合一致しますので，ただ書き変わっただけだと思っていただいて問題ありません。}。ただ，切片$b$を$\beta_0$として右辺の一番前に持って来ました。というのも，そうするとベクトルで書く時にわかりやすいからです。

説明変数行列の左端に1を入れたベクトルを追加し，回帰係数$\beta$もセットにして，次のように表現します。
\[
	\bm{X} = \begin{pmatrix}
		1      & x_{11} & x_{12} & \cdots & x_{1m} \\
		1      & x_{21} & x_{22} & \cdots & x_{2m} \\
		\vdots & \vdots & \vdots & \vdots & \vdots \\
		1      & x_{n1} & x_{n2} & \cdots & x_{nm}
	\end{pmatrix},
	\beta=\begin{pmatrix}
		\beta_0 \\
		\beta_1 \\
		\vdots  \\
		\beta_m
	\end{pmatrix}
\]
とすると，重回帰分析の式は次のように簡単な表現に変わります。
\[ \bm{Y} = \bm{X \beta} + \bm{e} \]

なんと，説明変数が$m$個に増えたのに，式の形は(単)回帰分析のそれと同じではありませんか！

このように，行列表現をすると多くの変数を一気に扱い，表現できるのです。このため，多変量解析ではベクトルの表記が基本になります。サイズを気にせず一般的に表現できるからです。

実際にこれらの式を読む時は，行列のサイズをイメージしながら読むと良いでしょう。たとえば左辺の$\bm{Y}$はサイズ$n$のベクトルなのですから，右辺の$\bm{X \beta}$もサイズ$n$の縦ベクトルになるはずなのです。実際，$\bm{X}$は$n \times (m+1)$の行列で，$\beta$は$(m+1) \times 1$のベクトルですから，計算結果は$n \times 1$，つまりサイズ$n$の縦ベクトルです。

\section{デザイン行列}
ところで，心理統計と言えば平均値の差を見ることだ，という話はこれまで散々聞いてきたところかと思います。
心理学は要因計画を立て，標本の平均値差から母集団に話を一般化するために，推測統計の知見を使って，帰無仮説検定やベイズ推定法を駆使するというやつです。
この要因計画は実は線形モデルの一環であり，\keyterm{一般線形モデル}{いっぱんせんけいもでる}{General Linear Model}と呼ばれています。
これを行列で表現することを，ここでは少し考えてみたいと思います。

まずは回帰分析と要因計画は何が同じで何が違うのかを，はっきりさせましょう。
同じところは線形モデルであるというところ，違うところは，回帰分析は説明変数も従属変数も連続変数であるのに対し，要因計画では一般に説明変数が離散変数であること，でした。離散変数であるとは，言い方を変えると名義尺度水準の数字だということです。すなわち「統制群」か「実験群」か，という違いを表すのに，$0,1$と言った数字を割り振ったものです。これは別に$3$と$12523$，という数字を割り振ったと言ってもいいのです。だって名義尺度水準は，数字と対象が一対一対応していれば良いのですから。

ここでは数学的に話をしやすくするために，統制群を$0$，実験群を$1$とするとしましょう。
線形モデルという枠組みは一緒なのですから，従属変数$y_i$が説明変数$x_i$によって変わるわけですが，ここではこの$x$が$ \{ 0, 1 \}$，というわけです。線形モデルを要素ごとに表現すると次のようになります。

\[ y_i = \beta_0 + \beta_1 x_i + e_i \]

この時，$i$さんは統制群に割り振られていたとすると，$x_i=0$ですから，この式は次のようになります。

\[ y_i = \beta_0 + e_i \]

逆に，$i$さんが実験群に割り振られていたとしますと，$x_i= 1$ですから，この式は次のようになります。

\[ y_i = \beta_0 + \beta_1 + e_i \]

これを見るとわかるように，両群のベースライン$\beta_0$は同じで，そこに効果$\beta_1$が乗っかるかどうか，が興味の的になります。この式の右辺に$i$は誤差成分しかなく，誤差を抜きにすると従属変数は$\beta_1$だけ変化するはずだ，というところから，平均値の差を検証しましょうと言ってることになるのでした。

これも$x_i$が個人ごとに変わるベクトルだと考えると，行列表現では次のようになります。

\[ \bm{y} = \bm{X\beta} + \bm{e} \]

同じ形ですね！ただし，ここでベクトル$\bm{x}$は，その中身が$\bm{x} = ( 0,0,0,1,1,0,1,0,...)$のように実験群か統制群かを分けるフラグが入っているだけになります。

以上は実験群と統制群という2群の話でしたが，3群以上になっても基本的なアイデアは同じです。
たとえば表\ref{tbl::07_01}のようなデータセットがあったとしましょう。
\begin{table}[ht]
	\centering
	\caption{群間要因(3水準)のデータセット例}
	\begin{tabular}{ccc}
		参加者番号 & 群わけ & 従属変数 \\\hline
		1     & A   & 3    \\
		2     & A   & 3    \\
		3     & A   & 4    \\
		4     & A   & 4    \\
		5     & B   & 6    \\
		6     & B   & 7    \\
		7     & B   & 8    \\
		8     & B   & 9    \\
		9     & C   & 7    \\
		10    & C   & 6    \\
		11    & C   & 5    \\
		12    & C   & 4    \\\hline
	\end{tabular}
	\label{tbl::07_01}
\end{table}

ここで群ごとの効果を表現したいとすると，次のように書くことになります。
\begin{eqnarray*}
	y_1 &=& \beta_0 + \beta_1 + e_1 \\
	y_2 &=& \beta_0 + \beta_1 + e_2 \\
	y_3 &=& \beta_0 + \beta_1 + e_3 \\
	y_4 &=& \beta_0 + \beta_1 + e_4 \\
	y_5 &=& \beta_0 + \beta_2 + e_5 \\
	y_6 &=& \beta_0 + \beta_2 + e_6 \\
	y_7 &=& \beta_0 + \beta_2 + e_7 \\
	y_8 &=& \beta_0 + \beta_2 + e_8 \\
	y_9 &=& \beta_0 + \beta_3 + e_9 \\
	y_{10} &=& \beta_0 + \beta_3 + e_{10} \\
	y_{11} &=& \beta_0 + \beta_3 + e_{11} \\
	y_{12} &=& \beta_0 + \beta_3 + e_{12}
\end{eqnarray*}
添字の対応に注意しながらみてくださいね。群Aの効果は$\beta_1$，群Bの効果は$\beta_2$，群Cの効果は$\beta_3$になります。

この$\beta$それぞれを該当するところ(割り当てられた群)だけに対応させつつ，統一的表現形である$\bm{y}=\bm{X\beta}+\bm{e}$にするには，次のように書く必要があります。
\[
	\begin{pmatrix}
		3 \\
		3 \\
		4 \\
		4 \\
		6 \\
		7 \\
		8 \\
		9 \\
		7 \\
		6 \\
		5 \\
		4
	\end{pmatrix}
	=
	\begin{pmatrix}
		1 & 1 & 0 & 0 \\
		1 & 1 & 0 & 0 \\
		1 & 1 & 0 & 0 \\
		1 & 1 & 0 & 0 \\
		1 & 0 & 1 & 0 \\
		1 & 0 & 1 & 0 \\
		1 & 0 & 1 & 0 \\
		1 & 0 & 1 & 0 \\
		1 & 0 & 0 & 1 \\
		1 & 0 & 0 & 1 \\
		1 & 0 & 0 & 1 \\
		1 & 0 & 0 & 1
	\end{pmatrix}
	\begin{pmatrix}
		\beta_0 \\
		\beta_1 \\
		\beta_2 \\
		\beta_3
	\end{pmatrix}+
	\begin{pmatrix}
		e_1    \\
		e_2    \\
		e_3    \\
		e_4    \\
		e_5    \\
		e_6    \\
		e_7    \\
		e_8    \\
		e_9    \\
		e_{10} \\
		e_{11} \\
		e_{12}
	\end{pmatrix}
\]
このような表記になった時の$\bm{X}$のことをとくに，実験のデザインを表している行列ということで，\keyterm{デザイン行列}{でざいんぎょうれつ}{design matrix}といいます。デザイン行列は自分で書くと面倒な感じがしますが，ともかくこのような書き方で$\bm{y}=\bm{X\beta}+\bm{e}$の統一表現は可能です。

ところでこのデザイン行列，$\bm{X}$のサイズは今回$n \times (m+1)$になっていますね($m$は水準数)。2水準のときは2列で済んだものが，3水準になると4列になるのはおかしくないですか？
そうです，1つ大事なポイントを忘れていました。各群の値はベースライン$\beta_0$からの相対的な違いです。相対的な違いというのは，言い換えると$\sum \beta =0$，すなわち全部の水準の和が$0$である必要があるのです。この式は今回の例だと$\beta_1+\beta_2+\beta_3=0$であり，これを移項すると明らかなように$\beta_3 = 0 - \beta_1 - \beta_2$です。つまり総和が決まっているので，自由に大きさを推定できるのは水準数$-1$になります。

これを踏まえてデザイン行列を次のように書き換えることができます。

\[
	\begin{pmatrix}
		3 \\
		3 \\
		4 \\
		4 \\
		6 \\
		7 \\
		8 \\
		9 \\
		7 \\
		6 \\
		5 \\
		4
	\end{pmatrix}
	=
	\begin{pmatrix}
		1 & 1  & 0  \\
		1 & 1  & 0  \\
		1 & 1  & 0  \\
		1 & 1  & 0  \\
		1 & 0  & 1  \\
		1 & 0  & 1  \\
		1 & 0  & 1  \\
		1 & 0  & 1  \\
		1 & -1 & -1 \\
		1 & -1 & -1 \\
		1 & -1 & -1 \\
		1 & -1 & -1
	\end{pmatrix}
	\begin{pmatrix}
		\beta_0 \\
		\beta_1 \\
		\beta_2 \\
	\end{pmatrix}+
	\begin{pmatrix}
		e_1    \\
		e_2    \\
		e_3    \\
		e_4    \\
		e_5    \\
		e_6    \\
		e_7    \\
		e_8    \\
		e_9    \\
		e_{10} \\
		e_{11} \\
		e_{12}
	\end{pmatrix}
\]

この式は先ほどの式と内容的には同じで，表現の仕方が違うだけですが，$\sum \beta= 0$の条件がなければ計算結果は算出されません。計算するための制約が少ないと，答えが出なくなるのです。$\sum \beta= 0$の制約条件を別途書き加えてもいいですが，制約条件も含めた行列の書き方ができるというわけですね。

線形モデルの統一的表現，あるいは行列の計算方法に少しは慣れてきたでしょうか。これがさらに多くの変数，未知数を扱うことになると，その利点はよりはっきりしてきます\footnote{ここでは触れませんが，被験者内計画・反復測定になっても行列表現はできます。個人差を表すデザイン行列を別途加えることになります。混合計画になると非常に複雑になりますが，それでも一般的な表現は可能です。}。

\section{ 因子分析モデルの行列表現}
ということで，因子分析モデルの代数的表現ですが，これも行列を使って表現すると非常にシンプルに表現できるということを説明していきましょう。

内容はまったく同じですが，確認しておきましょう。標準得点行列$\bm{Z}$を因子負荷行列$\bm{A}$と因子得点行列$\bm{F}$をつかって，次のように表します。
\begin{equation}
	\bm{Z} = \bm{F} \bm{A} ' + \bm{U} \bm{D}
\end{equation}

ここで，各行列の要素のサイズ感をつかんでおきましょう。まず$\bm{R}$というのは相関行列ですから，$m$個項目があるのでサイズは$m \times m$ の正方行列になります。次に$\bm{F}$ですが，これは因子得点の行列です。得点は人数分ありますから行は$n$，因子の数が列になるのでこれを$p$とすると$n\times p$です。$\bm{A}$は因子負荷行列。因子負荷行列は因子の数と項目の数の組み合わせだけあるわけですから，$m \times p$になりますね。$\bm{U}$は独自因子得点です。得点ですから人数分，独自性分は各項目にありますから，サイズとしては$n \times m$になります。最後に$\bm{D}$ですが，これは独自因子の負荷量です。項目の数だけあるのですが，列ベクトルや行ベクトルで表現すると計算の時にサイズが変わって不便なことになります。ですから，対角項に$d_j$をもつ正方行列$m \times m$として表現しています。

サイズを確認したところで，実際に行列計算をしてみましょう。

\begin{equation}
	\begin{array}{lll}
		\bm{R} & = & \frac{1}{N} \bm{Z}'\bm{Z}                                                                                                                                                      \\
		       &   & \hspace{8cm}\text{$\bm{Z}$を因子分析のモデル式にして}                                                                                                                                       \\
		       & = & \frac{1}{N}(\bm{F} \bm{A} ' + \bm{U} \bm{D})'(\bm{F} \bm{A} ' + \bm{U} \bm{D})                                                                                                 \\
		       &   & \hspace{8cm}\text{前の項の転置を中に入れます}                                                                                                                                               \\
		       & = & \frac{1}{N}\left \{ (\bm{F} \bm{A} ')'+(\bm{U} \bm{D})' \right \}(\bm{F} \bm{A} ' + \bm{U} \bm{D})                                                                             \\
		       &   & \hspace{8cm}\text{転置のカッコを外すときは順番を入れ替えて転置}                                                                                                                                      \\
		       & = & \frac{1}{N}(\bm{A} \bm{F} ' + \bm{D}' \bm{U} ')(\bm{F} \bm{A} ' + \bm{U} \bm{D})                                                                                               \\
		       &   & \hspace{8cm}                                                                                                                                                                   \\
		       & = & \frac{1}{N}\bm{A} \bm{F}' \bm{F} \bm{A}' + \frac{1}{N} \bm{A} \bm{F}' \bm{U} \bm{D} + \frac{1}{N} \bm{D}' \bm{U} ' \bm{F} \bm{A} ' + \frac{1}{N} \bm{D}' \bm{U}' \bm{U} \bm{D} \\
	\end{array}
\end{equation}
と，このように展開できました。記号を見ているとわかりにくいので，サイズ感を確認しましょう。最終的には，次のようになっています。
\[
	\underset{m \times m}{\bm{R}} =
	\frac{1}{N} \underset{m \times p}{\bm{A}} \underset{p \times n}{\bm{F}'} \underset{n \times p}{\bm{F}} \underset{p \times m}{\bm{A}'} +
	\frac{1}{N} \underset{m \times p}{\bm{A}} \underset{p \times n}{\bm{F}'} \underset{n \times m}{\bm{U}} \underset{m \times m}{\bm{D}} +
	\frac{1}{N} \underset{m \times m}{\bm{D}'} \underset{m \times n}{\bm{U}'} \underset{n \times p}{\bm{F}} \underset{p \times m}{\bm{A}'} +
	\frac{1}{N} \underset{m \times m}{\bm{D}'} \underset{m \times n}{\bm{U}'} \underset{n \times m}{\bm{U}} \underset{m \times m}{\bm{D}}
\]

ここで，要素ごとに計算していた時のことを思い出してください。第二項$\frac{1}{N} \bm{A} \bm{F}'\bm{U} \bm{D}'$と第三項$\bm{D}' \bm{U}' \bm{F} \bm{A}'$の中にある，$\bm{F}' \bm{U}$と$\bm{U}' \bm{F}$のところは，共通因子得点と独自因子得点の積ですし，いずれも標準化されていますから，$\frac{1}{N}$と合わせて考えると，これは相関係数を表していることになります。また，共通因子と独自因子は相関しませんので，これはイコール$\bm{0}$となり，この2つの項が消えてしまうのでした。

また，第一項の$\frac{1}{N}\bm{F} ' \bm{F}$は，共通因子同士の相関を表しています。$\bm{F'F}=\bm{C}$とすると，これは因子得点間相関$\bm{C}$を表すことになります。これが直交であると仮定する，つまり他の因子と相関しないと考えると，$\bm{C} = \bm{I}$，つまり単位行列です。単位行列は計算に影響を与えませんから，$\bm{AF'FA'}=\bm{ACA'}=\bm{AIA'}=\bm{AA'}$となり，この式は簡単に次のように変形できます。

\begin{equation}
	\bm{R} = \bm{A} \bm{A} ' + \bm{D}^2
\end{equation}
先ほどの代数的展開を，そのまま行列で表現しただけですが，この方がシンプルに表現できていますね。この表現は，因子分析の第一定理と第二定理の両方を含んで一度に表せているのです。

いかがでしょうか。行列表現の便利さがわかっていただければ，と思います。しかしこれでもまだ謎は残りますね。我々が追っている謎は，行列からどのようにして因子負荷量を計算するのか，です。それを知るためには，もう1つ線形代数から明らかになる特徴を知らなければなりません。次回をどうぞお楽しみに。
\section{課題}
\paragraph{行列計算を確認しておこう}$ \bm{V} = ( \bm{I} -\frac{1}{n}\bm{1} \bm{1} ')\bm{X} $の要素を書き下してみよう。平均偏差行列ができているでしょうか。
\paragraph{行列計算を確認しておこう2}$ \bm{S},\bm{Z},\bm{R}$も，面倒でも要素レベルまで書き下してみよう。
\paragraph{行列計算を確認しておこう3}因子分析モデルの行列計算の結果出てくる，$\bm{R} = \bm{A} \bm{A} ' + \bm{D}^2$の要素を確認し，エレメントワイズで表現していた式との対応を確認しよう。
\end{document}
